% !TeX root = ../main.tex

\chapter{Literature Review}\label{chapter:literature-review}

This chapter situates the thesis within four threads of prior work: the CLM problem
space and its enterprise stakes (Section~\ref{sec:lit:clm}), the evolution of contract
analysis methods from rule-based pipelines to LLMs
(Sections~\ref{sec:lit:rulebased} and~\ref{sec:lit:llm}), multi-agent systems and
their reception in legal AI (Section~\ref{sec:lit:mas}), the conditions under which
multi-agent designs help or hurt (Section~\ref{sec:lit:when}), and evaluation
methodology in legal NLP (Section~\ref{sec:lit:eval}). The chapter closes with a
synthesis (Section~\ref{sec:lit:synthesis}) that positions the present work against the
gap identified in Chapter~1.

\section{The CLM Problem Space}\label{sec:lit:clm}

Contract Lifecycle Management denotes the end-to-end set of activities through which
an organisation authors, negotiates, executes, monitors, renews, and ultimately retires
its contractual instruments. Practitioners typically partition the lifecycle into a
request phase, an authoring phase, a negotiation and approval phase, an execution
phase, and an obligation-management phase, with audit and renewal looping back into
authoring. From an information-systems perspective, CLM is therefore a long-running,
multi-stakeholder workflow that sits at the intersection of legal, procurement, finance,
and operational functions.

The economic stakes are large. World Commerce \& Contracting estimates that
organisations forfeit value equivalent to roughly nine percent of annual revenue through
poor contracting practices on average, with the figure rising substantially in complex
industries~\cite{worldcc2025benchmark}. The same source reports that contract-related
data is dispersed across an average of twenty-four systems within a single organisation,
and that the majority of executives surveyed regard their contracts as too rigid to
accommodate change~\cite{worldcc2025benchmark}. Goldman Sachs analysts estimate that
forty-four percent of legal-task time in the United States is technically exposed to
automation by generative AI, second only to administrative
occupations~\cite{goldmansachs2023ai}.

The bottleneck within this lifecycle is typically the review stage, in which contracts
must be parsed against playbooks to identify present clauses, missing clauses,
deviations from preferred wording, and risk-bearing provisions. Three sub-tasks
dominate. \emph{Clause extraction} locates and bounds the textual span that instantiates
a given clause category --- for example, an indemnification clause or a
change-of-control provision. \emph{Risk identification} classifies extracted clauses
against playbook standards, flagging deviations that warrant escalation.
\emph{Obligation tracking} converts executed clauses into a structured calendar of
duties, milestones, and contingent triggers. Each sub-task amplifies the cost of error:
a missed indemnification clause produces audit-time financial exposure, a misclassified
risk produces negotiation-time leverage loss, and an untracked obligation produces
post-execution non-compliance.

These three sub-tasks share a structural property that motivates the rest of this
chapter. They are all, at their core, span-extraction and span-classification problems
over long, heterogeneous, domain-specific documents --- exactly the regime in which
information-extraction methodology has evolved over the past two decades.

\section{From Rule-Based to Neural Approaches}\label{sec:lit:rulebased}

Pre-LLM contract analysis was dominated by symbolic and statistical methods. Rule-based
extractors, often built around regular expressions, gazetteers, and hand-engineered
grammars, achieved high precision on narrow clause types but degraded sharply on
contracts whose drafting style departed from the templates encoded in the rules. The
transition to statistical sequence labelling, exemplified by Conditional Random
Fields~\cite{lafferty2001crf}, provided a principled probabilistic treatment of span
boundaries and substantially improved cross-document robustness. The CRF formulation was
widely adopted for legal information extraction in the 2000s and remained competitive
through the deep-learning transition.

The deep-learning era introduced neural sequence labellers that combined bidirectional
recurrent encoders with CRF output layers. Lample et al.\ demonstrated that
bidirectional LSTM encoders coupled with a CRF decoding layer outperformed prior
named-entity recognition systems without hand-engineered features, establishing the
BiLSTM-CRF architecture as a workhorse for span extraction in legal and biomedical
domains~\cite{lample2016neural}. The introduction of the Transformer encoder and
BERT~\cite{devlin2019bert} initiated a further shift, replacing recurrent encoders with
self-attention and enabling effective transfer from large general-domain corpora to
small legal corpora through fine-tuning.

Domain-specific pre-training emerged in response to the observation that legal text
exhibits lexical, syntactic, and discourse properties that differ systematically from
general-purpose web text. LEGAL-BERT, introduced by Chalkidis et al., investigated
three adaptation strategies: out-of-the-box use, additional pre-training on legal
corpora, and pre-training from scratch, and reported that domain pre-training generally
improved downstream legal NLP performance, with pre-training from scratch giving the
largest gains in the most domain-shifted tasks~\cite{chalkidis2020legalbert}.
LEGAL-BERT and its descendants became standard baselines for legal classification and
extraction tasks throughout the late 2010s and early 2020s.

Three limitations of the rule-based and BERT-era stack motivated the transition to
LLMs. First, span-extraction supervision is expensive: legal annotation requires
lawyer-time, and the expert-annotated CUAD dataset, with its five hundred contracts and
forty-one categories~\cite{hendrycks2021cuad}, was a large undertaking by the standards
of legal NLP. Second, BERT-class encoders are limited by their context windows,
typically 512 tokens, which forces sliding-window inference over long contracts and
complicates document-level reasoning. Third, BERT-class models do not generalise
zero-shot to new clause categories: each new category requires labelled examples, which
is incompatible with the long tail of clause types that arises in production CLM. These
limitations created the opening that LLM-based contract analysis has, in the past four
years, attempted to fill.

\section{LLMs in Contract Analysis}\label{sec:lit:llm}

The CUAD dataset of Hendrycks et al.\ is the foundational benchmark for clause-level
contract analysis. It comprises 510 commercial contracts sourced from EDGAR filings,
with 13{,}000 expert annotations across forty-one clause
categories~\cite{hendrycks2021cuad}. The original baselines, fine-tuned RoBERTa- and
DeBERTa-class encoders, achieved best results in the regime of forty-four percent
precision at eighty percent recall (AUPR-based reporting), which the authors interpreted
as evidence of a substantial capability gap in 2021. CUAD remains the most widely cited
public benchmark for contract clause extraction and underpins the evaluation in this
thesis.

Two complementary datasets have shaped subsequent work. ContractNLI, introduced by
Koreeda and Manning, recasts a subset of contract review as document-level natural
language inference: given a non-disclosure agreement and a hypothesis drawn from a fixed
set of seventeen, the system must label entailment, contradiction, or neutrality, while
additionally identifying evidence spans~\cite{koreeda2021contractnli}. MAUD, introduced
by Wang et al., focuses on merger agreements and provides reading-comprehension
supervision derived from the American Bar Association's Public Target Deal Points Study,
with 150 agreements and roughly 39{,}000 examples~\cite{wang2023maud}. LEDGAR,
introduced by Tuggener et al.~\cite{tuggener2020ledgar}, scales further along the
multi-label classification axis with approximately 100{,}000 provisions extracted from
60{,}000 contracts, at the cost of weaker supervision quality. These datasets together
provide the empirical surface against which LLM-based contract analysis is currently
evaluated.

The arrival of GPT-4-class models prompted a wave of human-comparison studies. The
most-cited example is Martin et al.'s ``Better Call GPT'' study, which compared
frontier LLMs against junior lawyers and legal-process outsourcers on issue
identification and issue location in commercial contracts~\cite{martin2024bettercallgpt}.
The study reports that GPT-4-1106 attained an issue-determination F-score of 0.87,
tied with legal-process outsourcers and within one point of junior lawyers, while
completing review in under five minutes per document against the human comparators'
forty-three to fifty-six minutes. The study additionally reported a 99.97\% cost
reduction. Although the study is vendor-funded and its quantitative claims should
therefore be received critically, the directional finding --- that frontier LLMs are
competitive with junior human reviewers on standardised contract-review tasks --- is
consistent with subsequent independent evidence.

The most recent and methodologically careful evaluation of LLMs on CUAD is
ContractEval, by Liu et al.~\cite{liu2025contracteval}. ContractEval evaluates nineteen
LLMs (four proprietary and fifteen open-source) on clause-level legal-risk
identification across forty-one categories, sampling 4{,}128 data points from 102 CUAD
contracts. The paper reports four findings of direct relevance to the present thesis.
First, proprietary models outperform open-source models on both correctness and
effectiveness metrics. Second, scaling open-source models yields diminishing returns.
Third, ``thinking'' or reasoning-mode prompting improves output effectiveness ---
measured by Jaccard similarity between predicted and reference spans --- but degrades
correctness, suggesting that extended reasoning over-complicates simple extraction
decisions. Fourth, open-source models exhibit a \emph{laziness} failure mode in which
they default to ``no related clause'' responses even when the relevant text is present
in the contract; ContractEval quantifies this with a laziness rate metric. The paper
additionally argues for $F_2$ alongside $F_1$, on the grounds that legal-risk
asymmetries make missed clauses costlier than spurious clauses.

ContractEval is the closest predecessor to the present thesis in scope and method, and
we adopt three of its design choices: $F_1$ and $F_2$ as primary correctness metrics,
Jaccard similarity as an effectiveness metric, and explicit accounting of the laziness
or wrong-rejection rate. We depart from ContractEval in two respects. First, we evaluate
\emph{architectures} rather than \emph{models alone}, fixing the family of LLMs and
varying the orchestration around them. Second, we sample at the natural class
distribution of CUAD (approximately thirty percent positive) rather than at the more
positive-heavy ratio that some evaluations use, which we argue exposes over-extraction
more honestly.

Three persistent limitations of LLM-based contract analysis remain salient regardless
of evaluation choice. The first is hallucination: LLMs occasionally produce extracted
spans that are not present in the source contract, a behaviour that is unacceptable in
legal review and has motivated retrieval-grounded
designs~\cite{cui2024chatlaw,raptopoulos2025pakton}. The second is prompt sensitivity:
minor, intent-preserving changes in the prompt can produce substantial changes in
output, as quantified by Chatterjee et al.'s POSIX prompt-sensitivity
index~\cite{chatterjee2024posix}. The third is the over-extraction / under-extraction
tradeoff: zero-shot models tuned to maximise recall on positive-heavy samples can
over-extract on natural distributions, while models tuned to avoid spurious extractions
can under-extract on rare categories. These three limitations frame the multi-agent
intervention.

\section{Multi-Agent Systems for Legal AI}\label{sec:lit:mas}

Multi-agent LLM systems decompose a task across multiple LLM invocations that play
differentiated roles and communicate through structured messages. The motivation,
common to general-purpose multi-agent frameworks such as
AutoGen~\cite{wu2023autogen}, MetaGPT~\cite{hong2024metagpt},
CAMEL~\cite{li2023camel}, and Mixture-of-Agents~\cite{wang2024moa}, rests on three
claims: that role specialisation reduces the per-step prompt complexity each agent must
handle, that decomposition exposes intermediate reasoning to inspection and
intervention, and that ensembling across agents recovers errors that any single agent
would commit. Reflexion~\cite{shinn2023reflexion} and Tree-of-Thoughts~\cite{yao2023tot}
add iterative-refinement and search-style structure on top of single agents, while
ReAct~\cite{yao2023react} interleaves reasoning with tool calls, and
Self-Consistency~\cite{wang2023selfconsistency} aggregates multiple stochastic
completions. Generative Agents demonstrate that long-running social simulations can be
sustained by LLM agents with memory and reflection~\cite{park2023generative}. Wei et
al.'s Chain-of-Thought prompting~\cite{wei2022cot} sits beneath all of these as the
within-agent reasoning primitive.

In the legal domain, three multi-agent systems are most directly relevant to this
thesis. ChatLaw, in its 2024 revision by Cui et al., reformulates the original
retrieval-augmented Chinese legal LLM~\cite{cui2024chatlaw} as a multi-agent
collaborative legal assistant with knowledge-graph-enhanced mixture-of-experts and
standardised operating procedures modelled on law-firm workflows; the paper reports a
7.73-point improvement over GPT-4 on Lawbench. PAKTON, by Raptopoulos et al., is a
fully open-source multi-agent framework for question answering in long legal agreements,
structured around three specialised agents: an \emph{Archivist} that handles parsing
and embedding with hierarchical document trees, a \emph{Researcher} that performs
hybrid graph-aware retrieval, and an \emph{Interrogator} that performs multi-step
iterative reasoning in a ReAct style~\cite{raptopoulos2025pakton}. PAKTON was accepted
as an oral presentation at EMNLP 2025 and is benchmarked on ContractNLI and
LegalBench-RAG. L-MARS, by Wang and Yuan, organises legal question answering as a
four-agent workflow comprising query decomposition, agentic search across BM25 and
case-law APIs, judge-style verification of jurisdictional and temporal validity, and
final summarisation~\cite{wang2025lmars}; the system reports a 38-point improvement
over zero-shot single-agent on its LegalSearchQA benchmark.

These three systems are valuable demonstrations that multi-agent legal architectures
can be assembled, but each evaluates a single architectural choice on a single benchmark
and does not ablate the multi-agent design itself. The empirical question of which
multi-agent pattern is preferable for clause-level contract analysis under a fixed model
and task remains underdetermined by this body of work.

The taxonomic axes along which these systems differ are nevertheless relevant for the
present thesis. Architectures vary in \emph{topology} (hierarchical with an orchestrator
versus peer with mutual messaging), in \emph{routing} (static pipelines versus dynamic
LLM-mediated routing versus dictionary-based routing), in \emph{specialisation
granularity} (one specialist per task type versus one specialist per clause category
versus one specialist per reasoning step), and in \emph{verification structure} (single
pass versus reflective revision versus multi-agent debate).
Table~\ref{tab:lit:mas} summarises a representative selection of multi-agent legal and
general systems along these axes.

\begin{table}[t]
  \centering
  \small
  \begin{tabular}{lllll}
    \toprule
    \textbf{System} & \textbf{Domain} & \textbf{Topology} & \textbf{Routing} &
    \textbf{Verification} \\
    \midrule
    ChatLaw v2~\cite{cui2024chatlaw}       & Legal QA (Chinese) & Hierarchical  & SOP-based     & MoE consensus    \\
    PAKTON~\cite{raptopoulos2025pakton}    & Legal QA/contracts & Pipeline      & Static        & ReAct iterative  \\
    L-MARS~\cite{wang2025lmars}            & Legal QA/search    & Pipeline      & Static        & Judge agent      \\
    AutoGen~\cite{wu2023autogen}           & General            & Peer          & Conversational& None (default)   \\
    MetaGPT~\cite{hong2024metagpt}         & Software dev.      & Hierarchical  & SOP-based     & Review agent     \\
    MoA~\cite{wang2024moa}                 & General            & Layered ensemble & Broadcast  & Aggregator       \\
    \bottomrule
  \end{tabular}
  \caption{Representative multi-agent systems by topology, routing, and verification
  structure.}
  \label{tab:lit:mas}
\end{table}

The reported gains in this body of work, while real, must be interpreted carefully.
ChatLaw v2's 7.73-point improvement is reported relative to GPT-4 on Lawbench, not
against a comparable multi-agent ablation; the 38-point improvement reported by L-MARS
is computed against zero-shot single-agent on a benchmark whose construction is itself
part of the contribution. Without architectural ablations, these gains conflate the
multi-agent design choice with retrieval, prompting, and benchmark-construction choices
made in the same systems.

\section{When Multi-Agent Helps and When It Does Not}\label{sec:lit:when}

A small but rapidly growing body of recent work has begun to interrogate the conditions
under which multi-agent designs improve on single-agent designs. Three contributions are
central to the framing of the present thesis.

Gao et al.\ ask, in a paper titled ``Single-agent or Multi-agent Systems? Why Not
Both?'', whether the empirical advantage of multi-agent systems persists as base-model
capability improves~\cite{gao2025single}. Their answer, supported by experiments on
BigCodeBench, SWE-Bench, and the MAST trace set, is that the advantage diminishes for
frontier models such as o3 and Gemini~2.5~Pro, whose improved long-context reasoning,
persistent memory, and tool-use absorb the workload that previously required
decomposition. The authors propose hybrid architectures that localise the error-prone
agent inside an otherwise single-agent design. The implication for the present thesis is
that whether multi-agent helps in CLM may depend on the strength of the underlying
model.

Cemri et al., in ``Why Do Multi-Agent LLM Systems Fail?'', construct the Multi-Agent
System Failure Taxonomy (MAST) by analysing more than 200 tasks and 1{,}600 traces
across seven multi-agent frameworks including AutoGen, MetaGPT, ChatDev, and
CAMEL~\cite{cemri2025why}. They identify fourteen failure modes grouped into three
categories: specification and system-design failures (such as disobedience to role
specification and rigid step repetition), inter-agent misalignment failures (such as
conversation-history loss and goal drift), and verification-and-termination failures.
The authors report that simple interventions --- clearer role specification, tighter
orchestration --- yield only modest gains, of the order of fourteen to fifteen percent,
suggesting that multi-agent failures arise from fundamental design rather than from
prompt-level tuning.

Bogavelli et al.\ introduce AgentArch, a benchmark that evaluates eighteen agent
configurations across six frontier LLMs on enterprise workflow tasks~\cite{bogavelli2025agentarch}.
They report two findings of direct relevance. First, no single agent architecture is
universally optimal across model families: smaller models can outperform larger models
when matched to an appropriate architecture, and the architecture that works best for
one family does not necessarily work best for another. Second, function-calling
prompting outperforms ReAct-style prompting in multi-agent settings, where ReAct
produces hallucinated tool calls. Absolute performance is also reported as low:
best-of-class models reach only 35.3\% on the harder enterprise task and 70.8\% on the
easier one.

Reading these three contributions together with Sanwal's Layered Chain-of-Thought
paper~\cite{sanwal2025layered}, which proposes external verification at each reasoning
layer to address the unverified-intermediate-step problem of vanilla CoT, and with the
failure-mode evidence from ContractEval~\cite{liu2025contracteval}, a coherent picture
emerges. Multi-agent designs are not a free lunch. They help when the task is
decomposable, when individual specialists can be made meaningfully better than a
generalist, when the orchestration overhead does not introduce its own failure modes,
and when the underlying model is weak enough that decomposition still offers a usable
lever. They hurt when these conditions do not hold. The empirical task of the present
thesis is to determine which of these conditions hold for clause-level contract analysis
on CUAD.

\section{Evaluation Methodology in Legal NLP}\label{sec:lit:eval}

Evaluation methodology in legal NLP, particularly for clause-level extraction, has not
converged on a single protocol, and the choices made in published evaluations materially
affect the rankings they produce. We treat four methodological axes here: metric
selection, statistical testing, prompt-sensitivity treatment, and sampling design.

\paragraph{Metric selection.}
Token-level $F_1$ remains the default for span extraction, but the legal-risk asymmetry
between false negatives (missed clauses, with downstream financial and compliance
exposure) and false positives (spurious clauses, recoverable at human review) has
motivated the use of $F_\beta$ with $\beta > 1$. ContractEval explicitly adopts $F_2$
alongside $F_1$ on this rationale and additionally reports Jaccard similarity as a
measure of span effectiveness~\cite{liu2025contracteval}. We follow this convention.
The choice is not innocuous: $F_2$ favours architectures that retrieve aggressively, and
an architecture that ranks first on $F_2$ may rank last on precision.

\paragraph{Statistical testing.}
Dror et al.'s ``Hitchhiker's Guide to Testing Statistical Significance in Natural
Language Processing''~\cite{dror2018hitchhikers} remains the canonical reference for
significance testing in NLP. The guide argues that paired tests should be used when
systems are evaluated on the same instances, that test selection should reflect
distributional assumptions, and that multiple-comparison correction should be applied
when many systems are compared. For paired binary outcomes --- system A correct, system
B incorrect, and the reverse --- McNemar's test~\cite{mcnemar1947note} is the standard.
For continuous metrics on small samples, paired bootstrap is preferred. The present
thesis applies McNemar for binary correctness and paired bootstrap for $F_2$
differences, with Holm-style correction for the family of comparisons under each
hypothesis.

\paragraph{Prompt sensitivity.}
Chatterjee et al.'s POSIX index~\cite{chatterjee2024posix} formalises the observation
that LLM outputs change under intent-preserving prompt perturbations such as
paraphrasing, formatting changes, and template substitutions. They report that increased
parameter count and instruction tuning do not necessarily reduce sensitivity, and that a
single few-shot exemplar typically reduces sensitivity substantially. The implication
for legal evaluation is that single-prompt comparisons may be unstable, and the
architectural ranking between systems can swap under prompt perturbation. We treat
prompt sensitivity as a methodological constraint: we fix prompt templates per
architecture across models, document the prompts in full, and discuss prompt sensitivity
as a threat to validity in Chapter~6.

\paragraph{Sampling design.}
The most consequential methodological choice in CUAD-based evaluation is the sampling
ratio between clause-present and clause-absent instances. CUAD's natural distribution
at the clause-category level is heavily imbalanced toward absences. Evaluations that
resample to a more positive-heavy ratio --- for example, two-thirds positive --- produce
optimistic precision estimates because the prior probability of a clause being present
is artificially elevated. Evaluations that hold the ratio at the natural level expose
over-extraction more honestly but risk small-sample variance on rare categories. We
adopt the natural-distribution sampling design at approximately thirty percent positive,
holding the same sample across architectures and models so that comparisons are paired
at the instance level.

The combination of metric, statistical, sensitivity, and sampling choices defines a
protocol space, and protocols differ substantially across published evaluations. The
present thesis fixes a single protocol and applies it uniformly across all eleven models
and five architectures, and reports the protocol explicitly so that other workers can
replicate or contest it.

\section{Synthesis and Research Positioning}\label{sec:lit:synthesis}

The literature surveyed in this chapter establishes four facts that motivate the
specific design of the present thesis. First, contract analysis is an enterprise-critical
task with a substantial economic surface, an established benchmark in
CUAD~\cite{hendrycks2021cuad}, and a recent baseline-establishing evaluation in
ContractEval~\cite{liu2025contracteval}. Second, multi-agent systems for legal AI exist
and are advocated, with notable systems such as ChatLaw~\cite{cui2024chatlaw},
PAKTON~\cite{raptopoulos2025pakton}, and L-MARS~\cite{wang2025lmars}, but each
evaluates a single architecture and does not ablate the multi-agent design choice.
Third, recent cross-domain evidence --- Gao et al.~\cite{gao2025single}, Cemri et
al.~\cite{cemri2025why}, and Bogavelli et al.~\cite{bogavelli2025agentarch} ---
challenges the assumption that multi-agent designs are universally beneficial, and
identifies model-architecture interaction, fourteen named failure modes, and a
diminishing-returns regime as serious caveats. Fourth, evaluation methodology in legal
NLP is heterogeneous, and the choice of metric, statistical test, and sampling ratio
materially shapes the conclusions one can draw.

The gap that this thesis addresses is the absence of a rigorous, paired, statistically
tested comparison of multi-agent architectural patterns for clause-level contract
analysis on CUAD, conducted across multiple frontier-model families, under a natural
class distribution, and with explicit mechanistic diagnosis of failure modes. Sub-RQ1
asks which architectural pattern is most effective; Sub-RQ2 asks how multi-agent
compares to single-agent on accuracy and cost; Sub-RQ3 asks which architectural
decisions enable auditable reasoning traces under the EU AI Act~\cite{euaiact2024} and
GDPR~\cite{gdpr2016} regimes. The hypotheses H1, H2, and H3, framed in Chapter~1,
operationalise the comparison.

The methodology developed in Chapter~3 instantiates this protocol: five architectural
configurations spanning zero-shot, chain-of-thought, pipeline decomposition, LLM-routed
specialists, and consolidated specialist prompts; eleven frontier models from three
families; the CUAD benchmark restricted to long contracts; natural class-distribution
sampling at approximately thirty percent positive; $F_1$, $F_2$, and Jaccard
correctness and effectiveness metrics; and McNemar and paired-bootstrap statistical
tests with multiple-comparison correction. Chapter~4 specifies the implementation,
including the LangGraph-based orchestration~\cite{langgraph2024} and Langfuse-based
observability layer~\cite{langfuse2024} that secure the auditable-trace property
addressed by Sub-RQ3. Chapter~5 reports the results, which, to anticipate, do not
uniformly favour multi-agent designs and instead reveal a model-family-dependent
ranking, a classification-gate rejection cascade specific to pipeline decomposition, and
a near-complete equivalence between LLM-routed specialists and consolidated specialist
prompts. The interpretation of these findings, and their implications for enterprise CLM
SaaS design, is taken up in Chapters~6 and~7.
